% Part: first-order-logic
% Chapter: models-of-arithmetic
% Section: models-of-pa

\documentclass[../../../include/open-logic-section]{subfiles}

\begin{document}

\olfileid{mod}{mar}{mpa}
\section{د~$\Th{PA}$ مدلونه}

\begin{explain}
د~$\Th{TA}$ هر نامعياري مدل د~$\Th{PA}$ هم نامعياري مدل دے۔ موږ
پوهېږو چې د~$\Th{TA}$، او په پايله کښې د~$\Th{PA}$، نامعياري
مدلونه شته۔ دا هم پوهېږو چې داسې نامعياري مدلونه نامعياري «شمېرې»
لري، يعنې د تعريف ساحې داسې !!{element}s چې له ټولو معياري «شمېرو»
هاخوا دي۔ خو دا څنګه ترتيب شوي دي؟ څومره دي؟ ومو ليدل چې د کمزورې
تيورۍ~$\Th{Q}$ مدلونه حتا يوازې يوه نامعياري شمېره لرلے شي۔ خو دغه
ساده !!{structure}s د~$\Th{PA}$ يا~$\Th{TA}$ مدلونه نۀ دي۔

د~$\Th{PA}$ يا~$\Th{TA}$ د مدلونو د جوړښت د پوهېدو کلي دا ده چې
وګورو په دغو تيوريو کښې کوم حقيقتونه !!{derivable} دي۔ مثلاً
$\Th{PA}$ لا له مخکې
$\lforall[x][\eq/[x][x']]$ او
$\lforall[x][\lforall[y][\eq[(x+y)][(y+x)]]]$ ثابتوي، نو هغه ساده
جوړښتونه چې دا !!{sentence}s دروغوي د~$\Th{PA}$ د مدلونو په توګه
ردوي۔

فرض کړئ~$\Struct{M}$ د~$\Th{PA}$ مدل دے۔ که
$\Th{PA} \Proves !A$ وي، نو $\Sat{M}{!A}$۔ بيا د
$\Assign{\Obj{0}}{M}$ دپاره~$\nszero$، د~$\Assign{\prime}{M}$ دپاره
$\nssucc$، د~$\Assign{+}{M}$ دپاره~$\nsplus$، د
$\Assign{\times}{M}$ دپاره~$\nstimes$ او د~$\Assign{<}{M}$ دپاره
$\nsless$ کاروو۔ هره !!{sentence}~$!A$ بيا د~$\nszero$، $\nssucc$،
$\nsplus$، $\nstimes$ او~$\nsless$ په هکله کوم شرط بيانوي، او که
$\Sat{M}{!A}$ وي نو هغه شرط بايد پوره شي۔ مثلاً که
$\Sat{M}{!Q_1}$، يعنې
$\Sat{M}{\lforall[x][\lforall[y][(\eq[x'][y'] \lif \eq[x][y])]]}$،
نو $\nssucc$ بايد !!{injective} وي۔
\end{explain}

\begin{prop}
په~$\Struct{M}$ کښې $\nsless$ يو خطي سخت ترتيب دے، يعنې:
\begin{enumerate}
\item د هېڅ~$x \in \Domain{M}$ دپاره $x \nsless x$ نۀ وي۔
\item که $x \nsless y$ او $y \nsless z$ وي، نو $x \nsless z$۔
\item د هر $x\neq y$ دپاره يا $x\nsless y$، يا $y\nsless x$۔
\end{enumerate}
\end{prop}

\begin{proof}
$\Th{PA}$ دا ثابتوي:
\begin{enumerate}
\item $\lforall[x][\lnot x < x]$
\item $\lforall[x][\lforall[y][\lforall[z][((x < y \land y < z) \lif x < z)]]]$
\item $\lforall[x][\lforall[y][((x < y \lor y < x) \lor \eq[x][y]))]]$
\end{enumerate}
\end{proof}


\begin{prop}
\ollabel{prop:M-discrete} $\nszero$ د~$\Domain{M}$ تر ټولو کوچنے
!!{element} د~$\nsless$ په ترتيب کښې دے۔ د هر~$x$ دپاره
$x\nsless x^\nssucc$، او $x^\nssucc$ تر ټولو کوچنے
$\nsless$-لوي !!{element} دے۔ له صفر څخه پرته د هر~$x$ دپاره داسې يوازينے~$y$ شته چې
$y^\nssucc=x$ وي۔ (موږ~$y$ په~$\Struct{M}$ کښې
د~$x$ «مخکښېنی» بولو او په~$^\nssucc x$ ئې ښيو۔)
\textbf{د ژباړې د نسخې سمون (OLMOD-010):} سرچينې د يوازيني
مخکښيني دعوه د هر x دپاره کړې وه، چې صفر هم رانغاړي، حال دا چې د
Q2 بديهي اصل له مخې صفر د هېڅ څيز جانشين نۀ دے۔ دلته دعوه په سم
ډول له صفر څخه پرته غړو ته محدوده شوې ده۔
\end{prop}

\begin{proof}
تمرين۔
\end{proof}

\begin{prob}
په~$\Lang{L_A}$ کښې داسې !!{sentence}s ومومئ چې په~$\Th{PA}$ کښې
!!{derivable} وي (او ځکه په~$\Struct{N}$ کښې رښتونې وي) او په
\olref[mod][mar][mpa]{prop:M-discrete} کښې د~$\nszero$، $\nssucc$
او~$\nsless$ ځانګړتياوې تضمين کړي۔
\end{prob}

\begin{prop}
د~$\Struct{M}$ ټول معياري !!{element}s د~$\nsless$ له مخې له ټولو
نامعياري !!{element}s څخه کوچني دي۔
\end{prop}

\begin{proof}
$n$ د معياري !!{element}~$\Value{\num{n}}{M}$ د
$\Struct{M}$ غړي د لنډون په توګه کاروو۔ لا~$\Th{Q}$ د هر~$n\in\Nat$ دپاره
$\lforall[x][(x < \num{n}' \lif (\eq[x][\num{0}] \lor
  \eq[x][\num{1}] \lor \dots \lor \eq[x][\num{n}]))]$ ثابتوي۔
د~$\nsless \nszero$ شرط پوره کوونکي !!{element}s نشته۔ نو که~$n$ معياري او~$x$
نامعياري وي، $x\nsless n$ نۀ شو لرلے۔ د تعريف له مخې نامعياري
عنصر د هېڅ~$n\in\Nat$ دپاره $\Value{\num{n}}{M}$ نۀ وي، نو
$x\neq n$ هم دے۔ څرنګه چې $\nsless$ خطي ترتيب دے، بايد
$n\nsless x$ وي۔
\end{proof}

\begin{prop}
د~$\Domain{M}$ هر نامعياري !!{element}~$x$ د لاندې فرعي سټ يو غړے
دے:
\[
\dots ^{\nssucc\nssucc\nssucc}x \nsless ^{\nssucc\nssucc}x \nsless
^{\nssucc}x \nsless x \nsless x^{\nssucc} \nsless x^{\nssucc\nssucc}
\nsless x^{\nssucc\nssucc\nssucc} \nsless \dots
\]
موږ دا فرعي سټ د~$x$ \emph{بلاک} بولو او په~$[x]$ ئې ليکو۔ نه تر
ټولو کوچنے او نه تر ټولو لوے !!{element} لري۔ دا د هغو
$y\in\Domain{M}$ د سټ په توګه ځانګړې کېدے شي چې د کوم معياري~$n$
دپاره $x\nsplus n=y$ يا $y\nsplus n=x$ وي۔
\end{prop}

\begin{proof}
داسې سټ~$[x]$ تل شته، ځکه د~$\Domain{M}$ هر !!{element}~$y$
يوازينے جانشين~$y^\nssucc$ لري او که دغه غړے نامعياري وي،
يوازينے مخکښېنی~$^\nssucc y$ لري۔ د پرله‌پسې !!{element}s $y$ او~$y^\nssucc$
دپاره $y\nsless y^\nssucc$، او $y^\nssucc$ د~$\nsless$ له مخې
د~$\Domain{M}$ تر ټولو کوچنے هغه !!{element} دے چې $y$ ترې
$\nsless$ وي۔ څرنګه چې تل
$^\nssucc y\nsless y$ او $y\nsless y^\nssucc$، نو $[x]$ نه تر ټولو
کوچنے او نه تر ټولو لوے !!{element} لري۔ که $y\in[x]$ وي، نو
$x\in[y]$، ځکه بيا يا $y^{\nssucc\dots\nssucc}=x$ يا
$x^{\nssucc\dots\nssucc}=y$۔ که
$y^{\nssucc\dots\nssucc}=x$ وي (له~$n$ ځله~$\nssucc$ سره)، نو
$y\nsplus n=x$ او برعکس، ځکه (که~$n$ د~$\prime$ نښو شمېر وي)
$\Th{PA} \Proves \lforall[x][\eq[x^{\prime\dots\prime}][(x +
    \num{n})]]$۔
\end{proof}

\begin{prop}
که $[x]\neq[y]$ او $x\nsless y$ وي، نو د هر $u\in[x]$ او
$v\in[y]$ دپاره $u\nsless v$۔
\end{prop}

\begin{proof}
پام وکړئ چې $\Th{PA} \Proves
\lforall[x][\lforall[y][(x < y \lif (x' < y \lor x' = y))]]$۔
نو که $u\nsless v$ وي او $[u]\neq[v]$، د هر~$n$ دپاره
$u\nsplus n^\nssucc\nsless v$ هم لرو۔

هر $u\in[x]$ د~$\nsless y$ شرط پوره کوي: د فرض له مخې $x\nsless y$۔ که
$u\nsless x$ وي، د انتقال له مخې $u\nsless y$۔ او که
$x\nsless u$ خو $u\in[x]$ وي، نو د کوم~$n$ دپاره
$u=x\nsplus n^\nssucc$؛ نو د هغه حقيقت له مخې چې اوس مو ثابت کړ
$u\nsless y$۔

اوس فرض کړئ $v\in[y]$ د~$\nsless y$ شرط پوره کوي، يعنې د کوم معياري~$m$
دپاره $v\nsplus m^\nssucc=y$۔ بيا $v\nsless x$ کېدے نۀ شي، ګنې
$y=v\nsplus m^\nssucc\nsless x$ به و۔ همدارنګه $x\neq v$، ګنې
$x\nsplus m^\nssucc=v\nsplus m^\nssucc=y$ او $[x]=[y]$ به وو۔
نو $x\nsless v$۔ بيا د هر~$n$ دپاره
$x\nsplus n^\nssucc\nsless v$ هم۔ له دې امله که $x\nsless u$ او
$u\in[x]$ وي، نو $u\nsless v$۔ که $u\nsless x$ وي، د انتقال له
مخې $u\nsless v$۔

په پای کښې که $y\nsless v$ وي، نو $u\nsless v$، ځکه لکه چې ومو
ښودل $u\nsless y$ او $y\nsless v$۔
\end{proof}

\begin{cor}
که $[x]\neq[y]$ وي، نو $[x]\cap[y]=\emptyset$۔
\end{cor}

\begin{proof}
فرض کړئ $z\in[x]$ او $x\nsless y$۔ بيا د هر $u\in[y]$ دپاره
$z\nsless u$۔ که $z\in[y]$ هم واے، $z\nsless z$ به مو لرلے۔
د~$y\nsless x$ حالت هم ورته دے۔
\end{proof}

\begin{explain}
دا معنا لري چې بلاکونه پخپله داسې ترتيبېدے شي چې د~$\nsless$ درناوی
وکړي: $[x]\nsless[y]$ هله او يوازې هله چې $x\nsless y$؛ يا په
هم‌ارزه توګه، چې د هر $u\in[x]$ او $v\in[y]$ دپاره
$u\nsless v$ وي۔ څرګنده ده چې معياري بلاک~$[0]$ تر ټولو کوچنے بلاک
دے۔ له هېڅ نامعياري بلاک سره تقاطع نۀ لري، او هېڅ دوه نامعياري
بلاکونه هم يو بل نۀ پرې کوي۔ په ځانګړي ډول، د پرله‌پسې جانشينونو يا
مخکښينو په اخيستو بل بلاک ته نۀ شئ رسېدے۔
\end{explain}

\begin{prop}
که $x$ او~$y$ نامعياري وي، نو $x\nsless x\nsplus y$ او
$x\nsplus y\notin[x]$۔
\end{prop}

\begin{proof}
که $y$ نامعياري وي، نو $y\neq\nszero$۔
$\Th{PA} \Proves \lforall[x][\lforall[y][(y \neq \Obj{0} \lif x < (x+y))]]$۔
\textbf{د ژباړې د نسخې سمون (OLMOD-011):} سرچينې د دې عمومي
حقيقت په ليکلو کښې y ازاد پرېښے و؛ دلته پر y هم کلي سور ورزيات
شوے، لکه دعوه او ورپسې استعمال چې غواړي۔
اوس فرض کړئ $x\nsplus y\in[x]$۔ څرنګه چې
$x\nsless x\nsplus y$، د کوم معياري عدد دپاره
$x\nsplus n^\nssucc=x\nsplus y$۔ خو د جمع د اختصار قانون
$\Th{PA} \Proves \lforall[x][\lforall[y][\lforall[z][(\eq[(x+y)][(x+z)] \lif y = z)]]]$
دے۔ له دې به د کوم معياري~$n$ دپاره $y=n^\nssucc$ پايله شي؛ خو
$y$ نامعياري فرض شوې ده۔
\end{proof}

\begin{prop}
تر ټولو کوچنے نامعياري بلاک نشته۔
\end{prop}

\begin{proof}
$\Th{PA} \Proves \lforall[x][\lexists[y][(\eq[(y+y)][x] \lor
      \eq[(y+y)'][x])]]$؛ يعنې هر~$x$ پر~$2$ وېشل کېږي (ښايي پاتې
شونې~$1$ ولري)۔ که $x$ نامعياري وي، $y$ هم نامعياري دے۔ د مخکينۍ
قضيې له مخې $y\nsless y\nsplus y$ او $y\nsplus y\notin[y]$۔
بيا $y\nsless(y\nsplus y)^\nssucc$ او
$(y\nsplus y)^\nssucc\notin[y]$ هم۔ خو
$x=y\nsplus y$ يا $x=(y\nsplus y)^\nssucc$، نو $y\nsless x$ او
$y\notin[x]$۔
\end{proof}

\begin{prop}
تر ټولو لوے بلاک نشته۔
\end{prop}

\begin{proof}
تمرين۔
\end{proof}

\begin{prob}
وښيئ چې د~$\Th{PA}$ په يوه نامعياري مدل کښې تر ټولو لوے بلاک
نشته۔
\end{prob}

\begin{prop}
\ollabel{prop:blocks-dense}
د بلاکونو ترتيب ګڼ دے۔ يعنې که $x\nsless y$ او $[x]\neq[y]$ وي،
نو داسې بلاک~$[z]$ شته چې له دواړو بېل او د هغو تر منځ وي۔
\end{prop}

\begin{proof}
فرض کړئ $x\nsless y$۔ لکه مخکې، $x\nsplus y$ پر دوه وېشل کېږي
(ښايي پاتې شونې ولري): داسې $z\in\Domain{M}$ شته چې يا
$x \nsplus y = z \nsplus z$ يا
$x \nsplus y = (z \nsplus z)^\nssucc$۔ عنصر~$z$ د~$x$ او~$y$
«اوسط» دے، او $x\nsless z$ او $z\nsless y$۔
\textbf{د ژباړې د نسخې سمون (OLMOD-012):} سرچينې د دې ثبوت په
دوو معادلو کښې ناتعريفه جمع کارولې وه؛ دلته د ټول څپرکي ټاکلې
نامعياري جمع بېرته ايښوول شوې ده۔
\end{proof}

\begin{prob}
د~\olref[mod][mar][mpa]{prop:blocks-dense} مفصل ثبوت وليکئ۔
د~$z$ د شتون د تضمين دپاره $\Th{PA}$ بايد کومه !!{sentence}
!!{derive} کړي؟ ولې $x\nsless z$ او $z\nsless y$، او ولې
$[x]\neq[z]$ او $[z]\neq[y]$ دي؟
\end{prob}

\begin{explain}
له دې امله، په هر شمېرېدونکي نامعياري مدل کښې نامعياري بلاکونه د
ناطقو په شان ترتيب لري: يو !!a{denumerable} ګڼ خطي ترتيب بې له
سرحدي ټکو جوړوي۔ ښودل کېدے شي چې هر دوه داسې !!{denumerable}
ترتيبونه آيزومورف دي۔ نو د رښتيني حساب د هر دوه !!{enumerable}
نامعياري مدلونو $\Struct{M}_1$ او~$\Struct{M_2}$ هغه راکمونې چې
يوازې~$<$ او~$=$ لرونکې ژبې ته کېږي، آيزومورفې دي۔ په رښتيا،
آيزومورفيزم~$h$ داسې تعريفېدے شي: د~$\Struct{M_1}$ او
$\Struct{M_2}$ معياري برخې د معياري مدل~$\Struct{N}$ سره، او ځکه له
يو بل سره، آيزومورفې دي۔ د نامعياري برخې بلاکونه پخپله د ناطقو په
شان ترتيب شوي او ځکه آيزومورف دي؛ د بلاکونو يو آيزومورفيزم د هر
بلاک دننه د خپل‌سرو عناصرو په سره برابرولو غځېدے شي، او بيا د
$\Struct{M_1}$ کښې د~$x$ د جانشين انځور د~$\Struct{M_2}$ کښې د
$x$ د انځور جانشين وټاکل شي۔ له دې دا \emph{نۀ} راځي چې
$\mathfrak{M}_1$ او~$\mathfrak{M}_2$ د حساب په بشپړه ژبه کښې
آيزومورف دي (آيزومورفيزم تل د يوې !!a{language} په نسبت وي)، ځکه پر
$\Domain{M_1}$ او~$\Domain{M_2}$ د جمع او ضرب د تعريفولو نامتجانسې
لارې شته۔ (دا د واټ د يوې مشهورې قضيې څخه هم راځي چې د يوې بشپړې
تيورۍ د شمېرېدونکو مدلونو شمېر 2 کېدے نۀ شي۔)
\textbf{د ژباړې د نسخې سمون (OLMOD-013):} سرچينې د بلاکونو
شمېرېدنه له هر نامعياري مدل څخه بې له شرطه پايله کړې وه؛ دا يوازې د شمېرېدونکي مدل دپاره لازمېږي۔ دلته دغه شرط څرګند شوے
او د ورپسې جملې له شمېرېدونکو مدلونو سره برابر شوے دے۔
\end{explain}
\end{document}
